\section{The \sys{} Prototype}
\label{sec:system}


The prototype is a checkpointed LangGraph state machine that reaches the
browser, perception and storage only through typed ports, each replaceable by an
in-memory fake (the experiments use the fakes). We describe the parts the
evaluation relies on.

\subsection{Orchestration, Capture and Persistence}

Exploration is a breadth-first search over (state, action) pairs. A plain Python
reference loop, kept as a parity oracle, has been ported to a LangGraph
\texttt{StateGraph} (\texttt{scan}, \texttt{reason}, \texttt{act},
\texttt{observe}). The checkpointed dictionary holds the frontier, the discovery
order, the skips and counters, but not the claim set, which lives in the graph
store; the two are separate stores, so resuming a checkpoint does not consult the
claims again (Sect.~\ref{sec:rq1}). Limits such as depth and state cap are
configuration, so a resumed run cannot silently change its rules. A skipped
action returns straight to \texttt{act}, since routing it through \texttt{scan}
would add a spurious entry to the discovery order, which is the evidence for
cycles.

A Playwright crawler records the DOM, the CDP accessibility tree and a screenshot;
candidate actions come from actionable AX roles and receive content-addressed
identifiers. The normaliser $N$ strips known volatile attributes
(\texttt{nonce}, \texttt{data-session}, \texttt{data-csrf} and similar),
hexadecimal-looking generated ids, prefixed tokens, runs of ten or more digits,
whitespace and volatile CDP keys. These patterns are broad on purpose, and
Sect.~\ref{sec:fp} shows what they merge. The claim set lives in Redis and the
graph, revisit edges included, in Neo4j. Neither store, nor the browser, is used in
the experiments of this paper.

\subsection{Perception}
\label{sec:perception}

Perception turns a capture into a \emph{Semantic UI Map}: elements with role,
accessible name, tags, selector, visibility and enabled state, plus the
capabilities the state exposes (for example \texttt{delete-user}). The
deterministic channel correlates DOM elements with AX nodes by browser identifier,
then by role plus name, then by unique name, and never guesses at ambiguous
matches. DOM-only elements are kept, since an element hidden from the AX tree can still be an exposure.

A VLM channel (a Gemini Flash model behind a rate-limited provider interface) reads the screenshot and AX tree to supply judgements the DOM cannot,
for example that a styled \texttt{div} looks like a delete button. By design it only
annotates elements the deterministic channel already found, adding tags and
capabilities; it can neither add nor remove an element. That limits its influence
but does not make the verdict immune to it: Eqs.~(\ref{eq:forb}) and~(\ref{eq:req})
test membership of a subject in $\Sig(v)$. Subjects that are identifiers or
accessible names come from the DOM and AX tree, but subjects that are tags or
capabilities can be created or hidden by a wrong annotation, in either direction,
which falls under Proposition~\ref{prop:asym}. The channel is present in the code base but not wired
into the orchestrator and \emph{not evaluated} here, so the division of labour is a
design hypothesis. Canvas-rendered UI is invisible to the deterministic channel.

\subsection{Negation Engine}
\label{sec:negation}

The engine implements Eqs.~(\ref{eq:forb}) and~(\ref{eq:req}) deterministically.
It checks that a clause applies to the run's role, then matches subjects against
capabilities (the stronger evidence) and then against element tags and accessible
names, exactly after whitespace and case normalisation; substring matching is
avoided because an over-eager matcher is itself an over-reporting bias.
Each violation carries its clause type, a severity fixed by clause type, a
rationale and evidence for a human reviewer. Each test application ships its pages,
an \texttt{app.json}, a structured \texttt{policy.json} beside the English policy,
and a \texttt{violations.json} of planted defects.
